\documentclass{situationist-preprint}

\title{Impeccable Freedom}
\msid{REL-14-a}
\msfamily{REL-14}
\msdiscipline{Philosophy of Religion}
\mstype{Principal manuscript}
\mskeywords{heaven; impeccability; sourcehood; alternative possibilities; Frankfurt cases; free will defence; character}

\begin{document}
\maketitle

\begin{abstract}
Traditional theology holds that the blessed in heaven cannot sin, and that they are free. The free will defence seems to make this combination unstable: if freedom requires the ability to do otherwise, the blessed are not free, and if it does not, God could have created free creatures who never sin, and the defence loses its point. We argue that the blessed are free in the sense of sourcehood, though not in the sense of alternative possibilities. Their good actions issue from reasons-responsive mechanisms that are their own, and the impossibility of their sinning is a consequence of what those mechanisms are and of what the blessed see, not of any constraint that bypasses them. Frankfurt cases show that the absence of alternatives does not by itself defeat responsibility; the case of the blessed is stronger, since the absence of alternatives is produced by the actual sequence. The blessed differ from creatures created impeccable by the history of their goodness: their character was formed by a history God did not settle. This answers the question why God did not create the blessed directly, in the way REL-02-a requires.
\end{abstract}

\section{The problem}

Augustine distinguished the freedom of the first human beings, who were able not to sin, from the freedom of the blessed, who will not be able to sin, and held the second to be the greater \citep[XXII.30]{AugustineCiv}. The tradition has kept both claims: the blessed are impeccable, and they are free. Two problems arise. The first is general: if freedom requires the ability to do otherwise, a being who cannot sin is not free in doing good. The second is specific to theodicy. The free will defence explains moral evil as a cost of creaturely freedom; but if freedom is compatible with impeccability, God could have created free creatures who cannot sin, and the cost was unnecessary \citep{Sennett1999,PawlTimpe2009}.

The most developed answer, due to \citet{Sennett1999} and \citet{PawlTimpe2009}, is that the freedom of the blessed is \emph{derivative}. The blessed cannot now do otherwise, but their character was formed by earlier free choices in which they could; their present actions are free because they derive from that earlier exercise of libertarian freedom. The answer meets both problems, but it requires libertarianism about the earlier choices, and it makes the freedom of heaven a borrowed freedom whose value depends on alternatives once open.

We give an answer that requires neither. On the account of ACT-01-a, an agent acts freely if her action issues from a reasons-responsive mechanism that is her own. On the account of ACT-02, Frankfurt cases show that responsibility does not require alternative possibilities. We argue that the blessed satisfy the first account and that their case is a stronger version of the second; and we argue, following the framework of REL-02-a, that what distinguishes the blessed from creatures created impeccable is not their freedom but the source of their goodness.

\section{Framework}

\begin{definition}[Freedom as sourcehood]\label{d:source}
An agent is \emph{free as a source} in performing an action if the action issues from a mechanism that is reasons-responsive and is the agent's own, in the sense of ACT-01-a.
\end{definition}

\begin{definition}[Freedom as alternatives]\label{d:alt}
An agent is \emph{free as having alternatives} in performing an action if, holding fixed the past and the laws, she could have performed some other action instead.
\end{definition}

\begin{definition}[Impeccability]\label{d:imp}
An agent is \emph{impeccable} if there is no situation she can be in in which she does wrong.
\end{definition}

\begin{assumption}[The beatific state]\label{a:beat}
The blessed have a character in which the love of the good is settled, and they see the good clearly: in every situation they can be in, every consideration that bears on what to do is apparent to them for what it is.
\end{assumption}

Assumption~\ref{a:beat} is the content of the tradition's description of the blessed as confirmed in grace and enjoying the vision of God, stated in terms of character and information. It does not say that God prevents them from sinning; it says what they are like and what they see.

\begin{principle}[Wrongdoing and reasons]\label{p:wrong}
To do wrong is to act against the balance of reasons that bear on one's situation. No situation provides sufficient reason to do wrong.
\end{principle}

\section{The argument}

\begin{lemma}\label{l:imp}
The blessed are impeccable.
\end{lemma}

\begin{proof}
A reasons-responsive mechanism, operating as it does, recognizes the reasons present and acts on them. By Assumption~\ref{a:beat}, every consideration in every situation the blessed can be in is apparent to them for what it is, and their love of the good is settled, so that nothing in them inclines against acting on what they recognize. In no such situation is there sufficient reason to do wrong (Principle~\ref{p:wrong}). The mechanism therefore issues in a wrong action in none of them, and by Definition~\ref{d:imp} the blessed are impeccable.
\end{proof}

\begin{lemma}\label{l:rr}
The impeccability of the blessed is consistent with their mechanisms' being reasons-responsive and their own.
\end{lemma}

\begin{proof}
Reasons-responsiveness is a counterfactual property: the mechanism would recognize sufficient reasons to act otherwise in a range of situations in which such reasons are present, and would act on them in some (ACT-01-a). The blessed can act otherwise in this sense: had the good required a different action, they would have recognized it and done it. What they cannot do is act wrongly, and Lemma~\ref{l:imp} derives this from the operation of the mechanism itself, not from anything that interferes with it. A mechanism that never errs because it responds perfectly to reasons is the limiting case of reasons-responsiveness, not a failure of it. Nor is the mechanism installed by another in a way that bypasses the agent's practical reasoning: as Section~\ref{s:history} argues, it was formed through that reasoning.
\end{proof}

\begin{theorem}\label{t:main}
The blessed are free as sources in their good actions, though not free as having alternatives to doing good.
\end{theorem}

\begin{proof}
By Lemma~\ref{l:rr} and Definition~\ref{d:source}, they are free as sources. By Lemma~\ref{l:imp}, there is no situation in which they do wrong; in particular, holding fixed their character and vision, they cannot perform any wrong action instead of the good one they perform, and so they lack alternatives to doing good in the sense of Definition~\ref{d:alt}.
\end{proof}

The relation to Frankfurt cases is instructive. In a Frankfurt case, an agent acts on her own while a counterfactual intervener stands ready to make her act in the same way if she shows signs of doing otherwise; the alternative is removed, but the removal plays no role in what actually happens \citep{Frankfurt1969}. ACT-02 concludes that such an agent is responsible. In the case of the blessed, the alternative of sinning is removed not by an idle intervener but by the actual sequence: by the agent's own character and vision, which are what produce her action. If an agent is responsible when her alternatives are removed by something external that does not operate, she is at least as clearly the source of her action when they are removed by her own reasons operating as they do. This is the structure of Luther's \enquote{I can do no other}, which \citet[ch.~6]{Dennett1984} used to argue that inability to do otherwise, when it flows from the agent's character, enhances rather than diminishes responsibility. \citet{Wolf1990} defended a related asymmetry: acting rightly for the right reasons is compatible with being unable to do otherwise, while acting wrongly is not.

\section{The history of the blessed}\label{s:history}

If impeccable freedom is possible, why did God not create the blessed directly, with a settled love of the good and a clear vision, and so avoid the evils of the history that leads to heaven? The derivative-freedom view answers that the blessed would then not be free. Our account cannot give this answer, since we have argued that impeccable freedom consists in sourcehood, which does not require a past with alternatives.

The answer is given by REL-02-a. Among the goods of a world is creatures' goodness that is not brought about by God: the state in which a creature acts well and God does not bring it about, in the sense of strong actualization, that she does. By the diagonal lemma of REL-02-a, God cannot bring about any state of this form. A God who created a creature with a settled good character would have brought that character about; whatever the creature's freedom as a source, its goodness would be the product of God's settling it, and the good of independence would be absent. The blessed, by contrast, came by their character through a history of their own practical reasoning: choices, failures, repentance and growth, which God did not settle. Their goodness has its source in them, not only in the thin sense of Definition~\ref{d:source}, which a directly created creature might also satisfy, but in the sense that the character from which their actions issue is the outcome of a history God did not bring about.

\begin{corollary}\label{c:direct}
The state of the blessed, including the independent goodness of their character, cannot be brought about by God directly; it is reachable only through a history that God did not settle.
\end{corollary}

\begin{proof}
The state of the blessed includes that their good character was not brought about by God. By the diagonal lemma of REL-02-a, God cannot bring about a state whose description includes that God does not bring it about.
\end{proof}

The view requires no libertarian premise. The history God did not settle may have been deterministic or not; what matters, as REL-02-a argues, is the relation between the creature's conduct and God's agency, not the modal profile of the creature's choices. Nor does it require that self-formation start from nothing: ACT-08 concludes that Strawson's argument against ultimate sourcehood \citep{GStrawson1994} fails because self-formation need not start from nothing, and the blessed formed their character from a nature they did not choose, as all agents do.

\section{Objections and replies}

\begin{objection}[Robots in heaven]
If the blessed cannot do otherwise than good, they are like perfectly programmed machines, and their good conduct is no more praiseworthy than a thermostat's regulation.
\end{objection}

\begin{reply}
A thermostat does not recognize reasons. The blessed act from a mechanism that recognizes every reason present and would act differently were the reasons different; their inability to sin is the inability of a perfectly reasons-responsive agent to act against what she sees. Their character was not programmed but formed through their own practical reasoning, as Section~\ref{s:history} argues. Praise follows sourcehood, and they are sources.
\end{reply}

\begin{objection}[The free will defence collapses]
If sourcehood without alternatives is freedom, God could have created free creatures who never sin, so evil is not the price of freedom \citep{Sennett1999}.
\end{objection}

\begin{reply}
We grant that God could have created creatures who are free as sources and never sin. REL-02-a does not claim otherwise: its defence does not rest on the value of freedom but on the value of creatures' goodness not brought about by God, which God cannot bring about by any means. Corollary~\ref{c:direct} applies this to heaven.
\end{reply}

\begin{objection}[Luck]
Whether a person comes to have the settled character of the blessed depends on the circumstances of her history, which she did not control. Her impeccable goodness is a matter of luck.
\end{objection}

\begin{reply}
Every character is formed in circumstances the agent did not control, and the blessed are no exception. Constitutive and circumstantial luck affect the formation of character in heaven as on earth; they do not show that the resulting actions do not issue from the agent's own mechanism. The companion manuscript REL-14-b examines moral luck in an impeccable world.
\end{reply}

\section{Conclusion}

The blessed cannot sin, and they are free. Their freedom is sourcehood: their actions issue from reasons-responsive mechanisms that are their own, and their inability to sin is the effect of those mechanisms operating with a settled love of the good and a clear sight of it. They lack alternatives to doing good, as Frankfurt's agents lack alternatives, but for a better reason. What God could not have given them directly is not their freedom but the independence of their goodness, which comes from a history God did not settle.

\begin{thebibliography}{99}\small
\bibitem[Augustine(Civ)]{AugustineCiv} Augustine. \emph{De civitate Dei}. Cited by book and chapter.
\bibitem[Dennett(1984)]{Dennett1984} Dennett, D.~C. (1984). \emph{Elbow Room: The Varieties of Free Will Worth Wanting}. MIT Press.
\bibitem[Frankfurt(1969)]{Frankfurt1969} Frankfurt, H.~G. (1969). Alternate possibilities and moral responsibility. \emph{The Journal of Philosophy}, 66(23), 829--839.
\bibitem[Pawl and Timpe(2009)]{PawlTimpe2009} Pawl, T. and Timpe, K. (2009). Incompatibilism, sin, and free will in heaven. \emph{Faith and Philosophy}, 26(4), 398--419.
\bibitem[Sennett(1999)]{Sennett1999} Sennett, J.~F. (1999). Is there freedom in heaven? \emph{Faith and Philosophy}, 16(1), 69--82.
\bibitem[Strawson(1994)]{GStrawson1994} Strawson, G. (1994). The impossibility of moral responsibility. \emph{Philosophical Studies}, 75(1--2), 5--24.
\bibitem[Wolf(1990)]{Wolf1990} Wolf, S. (1990). \emph{Freedom Within Reason}. Oxford University Press.
\end{thebibliography}

\end{document}
